\documentclass[11pt,letterpaper]{article}

\usepackage[spanish,es-noshorthands,es-tabla]{babel}
\usepackage{fontspec}
\setmainfont{Source Serif Pro}
\usepackage{microtype}
\usepackage{geometry}
\usepackage{enumitem}
\usepackage{xcolor}
\usepackage{csquotes}
\usepackage{float}
\usepackage{booktabs}
\usepackage{array}
\usepackage{tikz}
\usepackage{pgfplots}
\pgfplotsset{compat=1.18}
\usetikzlibrary{arrows.meta,positioning,fit,backgrounds,calc}
\usepackage{xurl}
\usepackage[hidelinks]{hyperref}
\usepackage[backend=biber,style=apa,sorting=nyt,maxcitenames=2]{biblatex}
\DeclareLanguageMapping{spanish}{spanish-apa}

\geometry{margin=2.3cm}
\hbadness=10000
\setlength{\emergencystretch}{1em}
\setlength{\parindent}{0pt}
\setlength{\parskip}{0.5em}
\setlist{nosep}
\addbibresource{referencias.bib}
\AtBeginBibliography{\sloppy\small}
\setlength{\bibitemsep}{0.3em}

% Colores: gris pizarra para el cajero automático (tecnología consolidada),
% verde para SINPE Móvil (tecnología innovadora), azul oscuro para rótulos.
\definecolor{cCaj}{RGB}{110,116,128}
\definecolor{cSM}{RGB}{36,122,82}
\definecolor{cEje}{RGB}{26,58,102}
\definecolor{cGris}{RGB}{120,116,110}

\tikzset{
  hito/.style={circle, draw=#1, fill=#1, inner sep=2.6pt},
  rotulo/.style={font=\footnotesize, align=center, text width=2.9cm,
    execute at begin node={\hyphenpenalty=10000\relax}},
  anio/.style={font=\small\bfseries, text=#1},
}

\begin{document}

% --- Portada -----------------------------------------------------------------
\begin{titlepage}
\centering
\vspace*{3.2cm}
{\large\scshape Gestión del Desarrollo Tecnológico:\\
Estrategias y Análisis de Portfolio\par}
\vspace{0.5cm}
{\normalsize Módulo 8. Foro 8.3: Conflictos entre tecnologías innovadoras y consolidadas\par}
\vspace{2.6cm}
{\color{cEje}\rule{0.55\textwidth}{0.8pt}\par}
\vspace{0.7cm}
{\LARGE\bfseries SINPE Móvil y el cajero automático\par}
\vspace{0.5cm}
{\large La tecnología consolidada no desapareció,\\
pero ya no rinde igual\par}
\vspace{0.7cm}
{\color{cEje}\rule{0.55\textwidth}{0.8pt}\par}
\vfill
{\large Carlos Luis Rojas Aragonés\par}
\vspace{0.3cm}
{\normalsize Chief Technology Officer Program\par}
\vspace{0.3cm}
{\normalsize 26 de octubre de 2026\par}
\vspace*{1.5cm}
\end{titlepage}

% --- Cuerpo ------------------------------------------------------------------
\section*{1. El caso}

En Costa Rica ya casi nadie pregunta \enquote{¿tiene efectivo?}, sino
\enquote{¿le hago un SINPE?}. SINPE Móvil, la plataforma de pagos inmediatos del
Banco Central de Costa Rica (BCCR), es un buen ejemplo de lo que plantea este
foro: una tecnología innovadora con una cuota de mercado creciente que no
reemplazó a las tecnologías consolidadas (el efectivo, el cajero automático, la
transferencia bancaria tradicional y la tarjeta), pero sí comprometió su
rendimiento. Todas siguen existiendo; ninguna rinde como antes.

\section*{2. De los cheques al celular: casi 30 años de infraestructura}

La innovación no salió de la nada. El 17 de abril de 1997, el BCCR puso en
marcha SINPE, el Sistema Nacional de Pagos Electrónicos \parencite{bccr_sinpe}.
Su primer servicio fue la compensación de cheques, que hasta entonces tardaban
en promedio 20 días en hacerse efectivos. Desde el año 2000 el sistema ofreció
transferencias inmediatas, hoy llamadas PIN, pensadas para montos altos
\parencite{araujo2024}. Para 2022, SINPE ya integraba 17 servicios y liquidaba
cerca de 3600 millones de dólares diarios \parencite{alvarado2022}.

Sobre esa base, el 7 de mayo de 2015 nació SINPE Móvil: el cliente asocia su
número de celular a una cuenta en colones y puede enviar dinero a cualquier
banco conociendo solo el número del destinatario, gratis hasta ₡100\,000
diarios. Al inicio, la curiosidad superó al uso: en las primeras seis semanas
se inscribieron 117\,287 usuarios, pero apenas se hicieron 1732 transacciones
\parencite{telesemana2015}. La Figura~\ref{fig:linea} resume los hitos.

\begin{figure}[H]
\centering
\begin{tikzpicture}[x=1.95cm]
  \draw[cGris, line width=1.2pt] (-0.35,0) -- (7.35,0);
  \foreach \i/\a/\t/\c/\p in {
      0/1997/{SINPE: compensación de cheques}/cEje/1,
      1/2000/{Transferencias inmediatas (hoy PIN)}/cEje/-1,
      2/2015/{SINPE Móvil y tarjetas con chip}/cSM/1,
      3/2016/{Cuentas simplificadas}/cEje/-1,
      4/2019/{IBAN en todas las cuentas}/cEje/1,
      5/2020/{Pandemia: despega la adopción}/cSM/-1,
      6/2025/{Código propio en la factura electrónica}/cCaj/1,
      7/2028/{QR y alias (meta del BCCR)}/cSM/-1} {
    \node[hito=\c] at (\i,0) {};
    \node[anio=\c] at (\i,{0.38*\p}) {\a};
    \ifnum\p>0
      \node[rotulo, anchor=south] at (\i,0.62) {\t};
    \else
      \node[rotulo, anchor=north] at (\i,-0.62) {\t};
    \fi
  }
\end{tikzpicture}
\caption{Hitos de SINPE y de las reformas que lo acompañaron. Elaboración propia
con base en \Textcite{araujo2024}, \Textcite{bccr2025fiscal} y
\Textcite{arias2026}.}
\label{fig:linea}
\end{figure}

\section*{3. Cinco años lentos y luego una explosión}

En 2019, SINPE Móvil registró apenas 6 millones de pagos. La pandemia lo cambió
todo: con los confinamientos, muchos hogares lo usaron por primera vez para
pagar comida y compras a domicilio, y a partir de ahí actuó el efecto red: entre
más personas lo usaban, más razones tenían los demás para adoptarlo
\parencite{araujo2024}. En 2023 hubo 506 millones de transacciones; en 2024,
648 millones \parencite{bccr2025usuarios}, y en 2025, 747,5 millones por
₡12,53 billones \parencite{rivero2026} (Figura~\ref{fig:anual}).

\begin{figure}[H]
\centering
\begin{tikzpicture}
\begin{axis}[
  ybar, bar width=30pt, width=11cm, height=6.2cm,
  /pgf/number format/assume math mode=true,
  enlarge x limits=0.18, ymin=0, ymax=850,
  symbolic x coords={2019,2023,2024,2025}, xtick=data,
  ylabel={Millones de transacciones}, ylabel style={font=\small, text=cEje},
  tick label style={font=\small}, axis line style={cGris}, tick style={cGris},
  ymajorgrids, grid style={cGris!25},
  nodes near coords, nodes near coords style={font=\small\bfseries},
]
  \addplot[draw=cSM, fill=cSM!80] coordinates {(2019,6) (2023,506) (2024,649) (2025,748)};
\end{axis}
\end{tikzpicture}
\caption{Transacciones anuales de SINPE Móvil. Elaboración propia con datos de
\Textcite{araujo2024}, \Textcite{bccr2025usuarios} y \Textcite{rivero2026}.}
\label{fig:anual}
\end{figure}

La penetración es igual de llamativa. En 2024, más del 76\,\% de los mayores de
15 años eran usuarios activos, cada uno con unas 180 transacciones al año, y
Costa Rica ocupaba el cuarto lugar mundial en pagos rápidos por habitante,
detrás de PromptPay en Tailandia, Pix en Brasil y Swish en Suecia. El valor
promedio de cada pago bajó a unos ₡16\,900, señal de que se usa para el gasto
diario \parencite{araujo2024}. A mayo de 2026 había 4,8 millones de números
inscritos, 1,29 por adulto \parencite{ramirez2026}, y entre enero y agosto de
2026 se acumulaban 446,1 millones de transacciones \parencite{arias2026}.

\section*{4. Cuatro tecnologías consolidadas que siguen, pero rinden menos}

\textbf{El efectivo y el cajero automático.} El efectivo es el medio de pago más
caro: en 2017 su manejo costó el equivalente al 1\,\% del PIB. Seguía reinando
por dos razones: funciona siempre, para todos, tenga o no cuenta bancaria, y
liquida al instante. SINPE Móvil le quitó justamente esa segunda ventaja. Los
cajeros pasaron de 155 millones de transacciones en 2019 a 104 millones en 2023,
mientras SINPE Móvil subía de 6 a 506 millones (Figura~\ref{fig:cajeros}). El
ahorro por usar menos efectivo se estima en un 0,5\,\% del PIB al año
\parencite{araujo2024}.

\begin{figure}[H]
\centering
\begin{tikzpicture}
\begin{axis}[
  ybar, bar width=26pt, width=10.5cm, height=6.2cm,
  /pgf/number format/assume math mode=true,
  enlarge x limits=0.45, ymin=0, ymax=600,
  symbolic x coords={2019,2023}, xtick=data,
  ylabel={Millones de transacciones}, ylabel style={font=\small, text=cEje},
  tick label style={font=\small}, axis line style={cGris}, tick style={cGris},
  ymajorgrids, grid style={cGris!25},
  nodes near coords, nodes near coords style={font=\small\bfseries},
  legend style={at={(0.03,0.97)}, anchor=north west, draw=none,
    fill=none, font=\small, legend cell align=left},
  legend image code/.code={\draw[#1, fill] (0cm,-0.1cm) rectangle (0.3cm,0.12cm);},
]
  \addplot[draw=cCaj, fill=cCaj!75] coordinates {(2019,155) (2023,104)};
  \addplot[draw=cSM, fill=cSM!80] coordinates {(2019,6) (2023,506)};
  \legend{Cajeros automáticos, SINPE Móvil}
\end{axis}
\end{tikzpicture}
\caption{Transacciones en cajeros automáticos y en SINPE Móvil en Costa Rica.
Elaboración propia con datos de \Textcite{araujo2024}.}
\label{fig:cajeros}
\end{figure}

\textbf{La transferencia bancaria tradicional.} SINPE Móvil ya concentra cerca
del 80\,\% de la cantidad de transferencias interbancarias, pero solo el 4\,\%
de su valor \parencite{araujo2024}. Los montos grandes siguen viajando por los
canales de siempre; lo que se perdió fue el volumen minorista.

\textbf{La tarjeta.} Es la que mejor resiste. En 2022 había más de 9,2 millones
de tarjetas emitidas, el 99\,\% sin contacto, y cerca de 200\,000 datáfonos
\parencite{alvarado2022}. Su punto débil es el costo para el comercio: el BCCR
fija un tope de comisión del 2\,\% por venta con tarjeta nacional, más hasta
₡12\,500 mensuales por terminal \parencite{madrigal2023}. Para un pequeño
comercio, recibir un SINPE Móvil gratis frente a pagar esa comisión es una
diferencia directa en su margen.

\textbf{El negocio de pagos de los bancos.} Los bancos dejaron de cobrar
comisiones por transferencias que ahora son gratuitas, pero el BIS encontró que
sus gastos no financieros bajaron con respecto a un escenario simulado en el que
SINPE Móvil no hubiera despegado. Además, cayó la concentración del mercado de
pagos \parencite{araujo2024}. Es decir, perdieron una fuente de ingresos, pero
ganaron eficiencia y enfrentan más competencia.

La Tabla~\ref{tab:resumen} resume el patrón: la tecnología consolidada no
muere, se repliega hacia los casos que la nueva no cubre.

\begin{table}[H]
\centering
\small
\renewcommand{\arraystretch}{1.25}
\begin{tabular}{>{\raggedright\arraybackslash}p{3.3cm}
                >{\raggedright\arraybackslash}p{5.6cm}
                >{\raggedright\arraybackslash}p{5.6cm}}
\toprule
\textbf{Tecnología consolidada} & \textbf{Lo que conserva} & \textbf{Lo que perdió} \\
\midrule
Efectivo y cajero & Personas sin cuenta o sin conexión estable & Un tercio de las operaciones en cajero (2019-2023) \\
Transferencia tradicional & El 96\,\% del valor transferido: los montos altos & Cerca del 80\,\% de la cantidad de transferencias \\
Tarjeta & Comercio en línea, crédito, turistas & Pagos pequeños entre personas y en el comercio pequeño \\
Bancos & La relación con el cliente y la cuenta & Comisiones por transferencia; menor concentración \\
\bottomrule
\end{tabular}
\caption{Cómo quedó cada tecnología consolidada. Elaboración propia con base en
\Textcite{araujo2024} y \Textcite{alvarado2022}.}
\label{tab:resumen}
\end{table}

\section*{5. Las tensiones de una adopción masiva}

El éxito trajo conflictos nuevos. Desde el 1 de setiembre de 2025, los comercios
deben identificar los pagos con SINPE Móvil con un código propio en la factura
electrónica. El BCCR aclaró que no es un impuesto, sino una medida de
trazabilidad fiscal \parencite{bccr2025fiscal}, pero el debate sobre si el
sistema facilita la informalidad sigue abierto \parencite{rivero2026}. Además,
el propio BIS advierte que hay personas con menos habilidades digitales o con
conexión inestable que seguirán prefiriendo el efectivo \parencite{araujo2024}.
Esa es, precisamente, la trinchera donde sobrevive la tecnología consolidada.

\section*{6. ¿Qué factores fomentan la adopción en distintos países?}

El volumen del BIS sobre pagos rápidos compara experiencias de América Latina y
Asia \parencite{aurazo2024}. De ahí, y del caso costarricense, salen seis
factores:

\begin{enumerate}
  \item \textbf{Un banco central que opera y además promueve.} En Costa Rica, el
    BCCR fue operador de la plataforma y catalizador del acuerdo con los bancos
    \parencite{aurazo2024}. En Brasil, con Pix, también del banco central, más
    del 90\,\% de los adultos envió o recibió un pago entre julio de 2023 y
    julio de 2024 \parencite{aurazo2024}.
  \item \textbf{Interoperabilidad.} Los bancos aceptaron que el BCCR unificara
    sus registros de teléfonos. El dinero sale directo de la cuenta, sin pasar
    por una billetera, a diferencia de los modelos de dinero móvil de las
    telefónicas en África \parencite{araujo2024}. En Perú, el banco central
    obligó a las billeteras Yape y Plin a interoperar \parencite{aurazo2024}.
  \item \textbf{Costo cero para el usuario.} Es común a SINPE Móvil y a UPI en
    India, donde además pesaron la facilidad de crear aplicaciones, la
    protección de datos y la alianza con el sector privado
    \parencite{cornelli2024}.
  \item \textbf{Reformas previas.} Tarjetas con chip y sin contacto (2015),
    cuentas simplificadas con solo la cédula (2016) e IBAN para todas las
    cuentas (2019). La bancarización de los mayores de 15 años subió del 77\,\%
    en 2018 al 90\,\% en 2024 \parencite{araujo2024}.
  \item \textbf{Un detonante externo.} La pandemia obligó a probar el pago
    digital, y el efecto red hizo el resto \parencite{araujo2024}.
  \item \textbf{Una infraestructura que ya existía.} SINPE Móvil se montó sobre
    casi dos décadas de SINPE. No hubo que construir desde cero la red entre
    bancos, solo una capa nueva encima.
\end{enumerate}

\section*{7. Lo que viene}

El BCCR prepara un SINPE Móvil con códigos QR (estáticos y con el monto
incluido) y con alias, para no tener que compartir el número de teléfono. La
idea es llevarlo al comercio físico, al comercio en línea y al pago de servicios
públicos, con 2028 como referencia. El número de teléfono y el SMS se
mantendrán, y las comisiones todavía se estudian a la luz de Brasil e India
\parencite{arias2026}. Si eso ocurre, el próximo terreno en disputa será el
datáfono.

Les dejo dos preguntas: cuando SINPE Móvil llegue a la caja del comercio, ¿la
tarjeta correrá la misma suerte que el cajero? Y en sus países, ¿quién está
impulsando los pagos inmediatos: el banco central o el sector privado?

\printbibliography[title={Referencias}]

\end{document}
